\chapter{2023年天津卷}

\section{单选题}

\begin{problem}
已知集合 \(U=\{1,2,3,4,5\}\)，\(A=\{1,3\}\)，\(B=\{1,2,4\}\)，则 \(\complement_U B\cup A=\)（\quad）

\choices
    {\(\{1,3,5\}\)}
    {\(\{1,3\}\)}
    {\(\{1,2,4\}\)}
    {\(\{1,2,4,5\}\)}

\end{problem}

\begin{answer}
A
\end{answer}

\begin{solution}
由
\(\complement_U B=\{3,5\}\)，
故
\(\complement_U B\cup A=\{3,5\}\cup \{1,3\}=\{1,3,5\}\).
故选 A.
\end{solution}

\begin{problem}
“\(a^2=b^2\)”是“\(a^2+b^2=2ab\)”的（\quad）

\choices
    {充分不必要条件}
    {必要不充分条件}
    {充分必要条件}
    {既不充分又不必要条件}

\end{problem}

\begin{answer}
B
\end{answer}

\begin{solution}
由
\(a^2+b^2=2ab\)
得
\((a-b)^2=0\)，
故 \(a=b\)，从而 \(a^2=b^2\)，必要性成立.

但 \(a^2=b^2\) 只说明 \(a=\pm b\). 当 \(a=-b\ne 0\) 时，
\(a^2+b^2\ne 2ab\)，
故充分性不成立.

所以为必要不充分条件. 故选 B.
\end{solution}

\begin{problem}
若 \(a=1.01^{0.5}\)，\(b=1.01^{0.6}\)，\(c=0.6^{0.5}\)，则 \(a,b,c\) 的大小关系为（\quad）

\choices
    {\(c>a>b\)}
    {\(c>b>a\)}
    {\(a>b>c\)}
    {\(b>a>c\)}

\end{problem}

\begin{answer}
D
\end{answer}

\begin{solution}
因为 \(1.01>1\)，故指数函数 \(1.01^x\) 单调递增，所以
\(b=1.01^{0.6}>a=1.01^{0.5}\).
又因为 \(0.6<1.01\)，且 \(x^{0.5}\) 在 \([0,+\infty)\) 上单调递增，所以
\(c=0.6^{0.5}<1.01^{0.5}=a\).
故
\(b>a>c\).
故选 D.
\end{solution}

\begin{problem}
函数 \(f(x)\) 的图象如下图所示，则 \(f(x)\) 的解析式可能为（\quad）

\bitmapfigure[max width=0.62\linewidth,max totalheight=4.8cm]{img/2023/tianjin/q4_fig1.png}

\choices
    {\(\dfrac{5(\e^x-\e^{-x})}{x^2+2}\)}
    {\(\dfrac{5\sin x}{x^2+1}\)}
    {\(\dfrac{5(\e^x+\e^{-x})}{x^2+2}\)}
    {\(\dfrac{5\cos x}{x^2+1}\)}

\end{problem}

\begin{answer}
D
\end{answer}

\begin{solution}
由图分析可知：函数图象关于 \(y\) 轴对称，其为偶函数，且 \(f(-2)=f(2)<0\).

由
\(\frac{5\sin(-x)}{(-x)^2+1}=-\frac{5\sin x}{x^2+1}\)
且定义域为 \(\R\)，即 B 中函数为奇函数，排除.

当 \(x>0\) 时，
\(\frac{5(\e^x-\e^{-x})}{x^2+2}>0\)，\(\frac{5(\e^x+\e^{-x})}{x^2+2}>0\)，
即 A、C 中函数在 \((0,+\infty)\) 上函数值为正，排除.

故选 D.
\end{solution}

\begin{problem}
已知函数 \(f(x)\) 的一条对称轴为直线 \(x=2\)，一个周期为 4，则 \(f(x)\) 的解析式可能为（\quad）

\choices
    {\(\sin \left( \dfrac{\pi}{2}x \right)\)}
    {\(\cos \left( \dfrac{\pi}{2}x \right)\)}
    {\(\sin \left( \dfrac{\pi}{4}x \right)\)}
    {\(\cos \left( \dfrac{\pi}{4}x \right)\)}

\end{problem}

\begin{answer}
B
\end{answer}

\begin{solution}
A、B 两式的周期均为
\(\frac{2\pi}{\frac{\pi}{2}}=4\)，
C、D 两式的周期均为 8，先排除 C、D.

再检验 \(x=2\)：
\(\sin \left( \frac{\pi}{2}\times 2 \right)=\sin \pi=0\)，
所以 \(x=2\) 是 \(\sin \left( \frac{\pi}{2}x \right)\) 图象的对称中心，不是对称轴.

而
\(\cos \left( \frac{\pi}{2}\times 2 \right)=\cos \pi=-1\)，
故 \(x=2\) 是 \(\cos \left( \frac{\pi}{2}x \right)\) 的一条对称轴.

故选 B.
\end{solution}

\begin{problem}
已知 \(\{a_n\}\) 为等比数列，\(S_n\) 为数列 \(\{a_n\}\) 的前 \(n\) 项和，\(a_{n+1}=2S_n+2\)，则 \(a_4=\)（\quad）

\choices
    {3}
    {18}
    {54}
    {152}

\end{problem}

\begin{answer}
C
\end{answer}

\begin{solution}
设首项为 \(a_1\)，公比为 \(q\).

当 \(n=1\) 时，
\(a_2=2a_1+2\)，
即
\(a_1q=2a_1+2\).

当 \(n=2\) 时，
\(a_3=2(a_1+a_2)+2\)，
即
\(a_1q^2=2(a_1+a_1q)+2\).

联立解得
\(a_1=2\)，\(q=3\).
故
\(a_4=a_1q^3=2\times 27=54\).
故选 C.
\end{solution}

\begin{problem}
调查某种群花萼长度和花瓣长度，所得数据如图所示，其中相关系数 \(r=0.8245\)，下列说法正确的是（\quad）

\begin{center}
\begin{tabular}{cc}
\bitmapinclude[max width=6.0cm,max totalheight=4.8cm]{img/2023/tianjin/q7_fig1.png}
&
\bitmapinclude[max width=5.0cm,max totalheight=4.8cm]{img/2023/tianjin/q7_fig2.png}
\end{tabular}
\end{center}

\choices
    {花瓣长度和花萼长度没有相关性}
    {花瓣长度和花萼长度呈现负相关}
    {花瓣长度和花萼长度呈现正相关}
    {若从样本中抽取一部分，则这部分的相关系数一定是 0.8245}

\end{problem}

\begin{answer}
C
\end{answer}

\begin{solution}
由散点图可知，样本点总体从左下向右上分布，因此花瓣长度和花萼长度呈正相关.

又 \(r=0.8245\) 只对应原样本数据，抽取部分样本后相关系数不一定仍为 0.8245.

故选 C.
\end{solution}

\begin{problem}
在三棱锥 \(P-ABC\) 中，线段 \(PC\) 上的点 \(M\) 满足 \(PM=\dfrac{1}{3}PC\)，线段 \(PB\) 上的点 \(N\) 满足 \(PN=\dfrac{2}{3}PB\)，则三棱锥 \(P-AMN\) 和三棱锥 \(P-ABC\) 的体积之比为（\quad）

\choices
    {\(\dfrac{1}{9}\)}
    {\(\dfrac{2}{9}\)}
    {\(\dfrac{1}{3}\)}
    {\(\dfrac{4}{9}\)}

\end{problem}

\begin{answer}
B
\end{answer}

\begin{solution}
如图，分别过 \(M,C\) 作 \(MM'\perp PA\)，\(CC'\perp PA\)，垂足分别为 \(M',C'\). 过 \(B\) 作 \(BB'\perp \text{平面 }PAC\)，垂足为 \(B'\)，连接 \(PB'\)，过 \(N\) 作 \(NN'\perp PB'\)，垂足为 \(N'\).

\bitmapfigure[max width=0.55\linewidth,max totalheight=4.8cm]{img/2023/tianjin/q8_solution_fig1.png}

因为 \(BB'\perp \text{平面 }PAC\)，\(BB'\subset \text{平面 }PBB'\)，所以 \(\text{平面 }PBB'\perp \text{平面 }PAC\).

又因为 \(\text{平面 }PBB'\cap \text{平面 }PAC=PB'\)，\(NN'\perp PB'\)，\(NN'\subset \text{平面 }PBB'\)，所以 \(NN'\perp \text{平面 }PAC\)，且 \(BB'//NN'\).

在 \(\triangle PCC'\) 中，因为 \(MM'\perp PA\)，\(CC'\perp PA\)，所以 \(MM'//CC'\)，所以
\(\frac{PM}{PC}=\frac{MM'}{CC'}=\frac{1}{3}\)
在 \(\triangle PBB'\) 中，因为 \(BB'//NN'\)，所以
\(\frac{PN}{PB}=\frac{NN'}{BB'}=\frac{2}{3}\)

所以
\[
\frac{V_{P-AMN}}{V_{P-ABC}}
=\frac{V_{N-PAM}}{V_{B-PAC}}
=\frac{\frac{1}{3}S_{\triangle PAM}\cdot NN'}{\frac{1}{3}S_{\triangle PAC}\cdot BB'}
=\frac{\frac{1}{3}\cdot \frac{1}{2}\left( PA\cdot MM' \right)\cdot NN'}{\frac{1}{3}\cdot \frac{1}{2}\left( PA\cdot CC' \right)\cdot BB'}
=\frac{2}{9}
\]
故选 B.
\end{solution}

\begin{problem}
双曲线 \(\frac{x^2}{a^2}-\frac{y^2}{b^2}=1\)（\(a>0\)，\(b>0\)） 的左、右焦点分别为 \(F_1,F_2\). 过 \(F_2\) 作其中一条渐近线的垂线，垂足为 \(P\). 已知 \(PF_2=2\)，直线 \(PF_1\) 的斜率为 \(\dfrac{\sqrt{2}}{4}\)，则双曲线的方程为（\quad）

\bitmapfigure[max width=0.42\linewidth,max totalheight=4.8cm]{img/2023/tianjin/q9_fig1.png}

\choices
    {\(\dfrac{x^2}{8}-\dfrac{y^2}{4}=1\)}
    {\(\dfrac{x^2}{4}-\dfrac{y^2}{8}=1\)}
    {\(\dfrac{x^2}{4}-\dfrac{y^2}{2}=1\)}
    {\(\dfrac{x^2}{2}-\dfrac{y^2}{4}=1\)}

\end{problem}

\begin{answer}
D
\end{answer}

\begin{solution}
设右焦点 \(F_2(c,0)\)，取渐近线
\(y=\frac{b}{a}x\).
则点到直线距离公式给出
\(PF_2=\frac{bc}{\sqrt{a^2+b^2}}=2\).
又双曲线满足 \(c^2=a^2+b^2\)，故
\(PF_2=\frac{bc}{c}=b=2\).

再由图可知
\(\tan \angle POF_2=\frac{b}{a}\)，
从而 \(OP=a\).

设 \(P\left( \frac{a^2}{c},\frac{ab}{c} \right)\)，由
\(k_{PF_1}=\frac{\frac{ab}{c}}{\frac{a^2}{c}+c}\)
并代入 \(b=2\)，\(c^2=a^2+4\)，得
\(\frac{a}{a^2+2}=\frac{\sqrt2}{4}\).
解得
\(a=\sqrt2\).
故所求双曲线方程为
\[
\frac{x^2}{2}-\frac{y^2}{4}=1
\]
故选 D.
\end{solution}
\section{填空题}

\begin{problem}
已知 \(\i\) 是虚数单位，化简 \(\frac{5+14\i}{2+3\i}\) 的结果为\fillinblank{}.

\end{problem}

\begin{answer}
\(4+\i\)
\end{answer}

\begin{solution}
有 \(\frac{5+14\i}{2+3\i}=\frac{(5+14\i)(2-3\i)}{(2+3\i)(2-3\i)}=\frac{52+13\i}{13}=4+\i\).
\end{solution}

\begin{problem}
在 \(\left( 2x^3-\frac{1}{x} \right)^6\) 的展开式中，\(x^2\) 项的系数为\fillinblank{}.

\end{problem}

\begin{answer}
60
\end{answer}

\begin{solution}
展开式通项为 \(T_{k+1}=\mathrm{C}_6^k (2x^3)^{6-k}\left(-\frac{1}{x}\right)^k=(-1)^k 2^{6-k}\mathrm{C}_6^k x^{18-4k}\). 令 \(18-4k=2\)，得 \(k=4\). 故系数为 \((-1)^4 2^2 \mathrm{C}_6^4=4\times 15=60\).
\end{solution}

\begin{problem}
过原点的一条直线与圆 \(C:(x+2)^2+y^2=3\) 相切，交曲线 \(y^2=2px\ (p>0)\) 于点 \(P\)，若 \(|OP|=8\)，则 \(p\) 的值为\fillinblank{}.

\end{problem}

\begin{answer}
6
\end{answer}

\begin{solution}
设切线方程为 \(y=kx\)（\(k>0\)）. 直线与圆相切，所以圆心 \((-2,0)\) 到直线 \(y=kx\) 的距离为 \(\sqrt3\)，即
\(\frac{2|k|}{\sqrt{1+k^2}}=\sqrt3\).
解得
\(k=\sqrt3\).

联立 \(\begin{cases} y=\sqrt3 x,\\ y^2=2px, \end{cases}\) 得除原点外
\(x=\frac{2p}{3}\)，\(y=\frac{2\sqrt3 p}{3}\).
故
\(|OP|=\sqrt{\left( \frac{2p}{3} \right)^2+\left( \frac{2\sqrt3 p}{3} \right)^2}=\frac{4p}{3}\).
由 \(|OP|=8\) 得 \(\frac{4p}{3}=8\)，解得 \(p=6\).
\end{solution}

\begin{problem}
甲、乙、丙三个盒子中装有一定数量的黑球和白球，其总数之比为 \(5:4:6\). 这三个盒子中黑球占总数的比例分别为 \(40\%,25\%,50\%\). 现从三个盒子中各取一个球，取到的三个球都是黑球的概率为\fillinblank{}；将三个盒子混合后任取一个球，是白球的概率为\fillinblank{}.

\end{problem}

\begin{answer}
0.05；\(\dfrac{3}{5}\)
\end{answer}

\begin{solution}
从三个盒子中各取一个球，均为黑球的概率为 \(\;0.4\times 0.25\times 0.5=0.05\).

设三个盒子球数分别为 \(5n,4n,6n\). 则黑球数分别为
\(2n,\quad n,\quad 3n\)，
白球总数为
\(15n-(2n+n+3n)=9n\).
故混合后任取一球为白球的概率为 \(\frac{9n}{15n}=\frac{3}{5}\).
\end{solution}

\begin{problem}
在 \(\triangle ABC\) 中，\(\angle A=60^\circ\)，\(BC=1\)，点 \(D\) 为 \(AB\) 的中点，点 \(E\) 为 \(CD\) 的中点. 若设 \(\overrightarrow{AB}=\bs{a}\)，\(\overrightarrow{AC}=\bs{b}\)，则 \(\overrightarrow{AE}\) 可用 \(\bs{a},\bs{b}\) 表示为\fillinblank{}；若 \(\overrightarrow{BF}=\dfrac{1}{3}\overrightarrow{BC}\)，则 \(\overrightarrow{AE}\cdot \overrightarrow{AF}\) 的最大值为\fillinblank{}.

\end{problem}

\begin{answer}
\(\dfrac{1}{4}\bs{a}+\dfrac{1}{2}\bs{b}\)；\(\dfrac{13}{24}\)
\end{answer}

\begin{solution}
\bitmapfigure[max width=0.42\linewidth,max totalheight=4.8cm]{img/2023/tianjin/q14_solution_fig1.png}

空 1：因为 \(D\) 为 \(AB\) 的中点，\(E\) 为 \(CD\) 的中点，所以 \(\overrightarrow{AD}=\frac{1}{2}\bs{a}\)，且
\(2\overrightarrow{AE}=\overrightarrow{AD}+\overrightarrow{AC}\).
故 \(\overrightarrow{AE}=\frac{1}{4}\bs{a}+\frac{1}{2}\bs{b}\).

空 2：由 \(\overrightarrow{BF}=\dfrac{1}{3}\overrightarrow{BC}\) 得 \(\overrightarrow{AF}=\frac{2}{3}\bs{a}+\frac{1}{3}\bs{b}\). 于是
\[
\overrightarrow{AE}\cdot \overrightarrow{AF}
=\left( \frac{1}{4}\bs{a}+\frac{1}{2}\bs{b} \right)\cdot \left( \frac{2}{3}\bs{a}+\frac{1}{3}\bs{b} \right)
=\frac{1}{12}\left( 2|\bs{a}|^2+5\bs{a}\cdot\bs{b}+2|\bs{b}|^2 \right)
\]

设 \(AB=x\)，\(AC=y\)，则 \(\bs{a}\cdot\bs{b}=xy\cos 60^\circ=\frac{xy}{2}\)，故
\(\overrightarrow{AE}\cdot \overrightarrow{AF}=\frac{1}{12}\left( 2x^2+\frac{5xy}{2}+2y^2 \right)\).
又由余弦定理，
\(BC^2=x^2+y^2-2xy\cos 60^\circ=x^2+y^2-xy=1\).
所以
\[
\overrightarrow{AE}\cdot \overrightarrow{AF}
=\frac{1}{12}\left( 2xy+\frac{5xy}{2}+2 \right)
=\frac{1}{12}\left( \frac{9xy}{2}+2 \right)
\]
由 \(x^2+y^2-xy=1\ge xy\) 得 \(xy\le 1\)，当且仅当 \(x=y=1\) 时取等号. 故最大值为
\(\frac{1}{12}\left( \frac{9}{2}+2 \right)=\frac{13}{24}\).
\end{solution}

\begin{problem}
若函数 \(f(x)=ax^2-2x-\left| x^2-ax+1 \right|\) 有且仅有两个零点，则 \(a\) 的取值范围为\fillinblank{}.

\end{problem}

\begin{answer}
\((-\infty,0)\cup (0,1)\cup (1,+\infty)\)
\end{answer}

\begin{solution}
分两种情况讨论.

当 \(x^2-ax+1\ge 0\) 时，\(f(x)=0\Longleftrightarrow (a-1)x^2+(a-2)x-1=0\Longleftrightarrow \bigl[(a-1)x-1\bigr](x+1)=0\).

若 \(a=1\)，则方程只有根 \(x=-1\)，且
\((-1)^2-a(-1)+1=3\ge 0\).

若 \(a\ne 1\)，则方程根为 \(x=-1\) 或 \(x=\dfrac{1}{a-1}\).

当根为 \(x=-1\) 时，需满足
\((-1)^2-a(-1)+1\ge 0\)，
即 \(a\ge -2\).

当根为 \(x=\dfrac{1}{a-1}\) 时，需满足
\(\left( \frac{1}{a-1} \right)^2-a\cdot \frac{1}{a-1}+1\ge 0\)，
解得 \(a\le 2\) 且 \(a\ne 1\).

当 \(x^2-ax+1<0\) 时，\(f(x)=0\Longleftrightarrow (a+1)x^2-(a+2)x+1=0\Longleftrightarrow \bigl[(a+1)x-1\bigr](x-1)=0\).

若 \(a=-1\)，则方程只有根 \(x=1\)，但
\(1-a+1=3>0\)，
不满足 \(x^2-ax+1<0\).

若 \(a\ne -1\)，则方程根为 \(x=1\) 或 \(x=\dfrac{1}{a+1}\).

当根为 \(x=1\) 时，需满足
\(1-a+1<0\)，
即 \(a>2\).

当根为 \(x=\dfrac{1}{a+1}\) 时，需满足
\(\left( \frac{1}{a+1} \right)^2-a\cdot \frac{1}{a+1}+1<0\)，
解得 \(a<-2\).

综上：
\[
\begin{aligned}
&a<-2 \text{ 时，零点为 } \frac{1}{a-1},\ \frac{1}{a+1};\\
&-2\le a<0 \text{ 时，零点为 } \frac{1}{a-1},\ -1;\\
&a=0 \text{ 时，只有一个零点 } -1;\\
&0<a<1 \text{ 时，零点为 } \frac{1}{a-1},\ -1;\\
&a=1 \text{ 时，只有一个零点 } -1;\\
&1<a\le 2 \text{ 时，零点为 } \frac{1}{a-1},\ -1;\\
&a>2 \text{ 时，零点为 } -1,\ 1.
\end{aligned}
\]

故 \(a\in (-\infty,0)\cup (0,1)\cup (1,+\infty)\).
\end{solution}
\section{解答题}

\begin{problem}
在 \(\triangle ABC\) 中，角 \(A,B,C\) 所对的边分别是 \(a,b,c\). 已知 \(a=\sqrt{39}\)，\(b=2\)，\(\angle A=120^\circ\).

\begin{enumerate}
\item 求 \(\sin B\) 的值；
\item 求 \(c\) 的值；
\item 求 \(\sin(B-C)\).
\end{enumerate}

\end{problem}

\begin{solution}
（1）由正弦定理，\(\frac{a}{\sin A}=\frac{b}{\sin B}\)，即 \(\frac{\sqrt{39}}{\sin 120^\circ}=\frac{2}{\sin B}\). 解得
\(\sin B=\frac{\sqrt{13}}{13}\).

（2）由余弦定理，\(a^2=b^2+c^2-2bc\cos A\). 代入已知得
\[
39=4+c^2-4c\left( -\frac{1}{2} \right)
\]
即 \(c^2+2c-35=0\). 解得 \(c=5\) 或 \(c=-7\)，舍去负值，故 \(c=5\).

（3）由正弦定理，\(\frac{a}{\sin A}=\frac{c}{\sin C}\)，得
\(\sin C=\frac{5\sqrt{13}}{26}\).
又 \(A=120^\circ\)，故 \(B,C\) 均为锐角，从而
\(\cos B=\sqrt{1-\sin^2 B}=\frac{2\sqrt{39}}{13}\)，\(\cos C=\sqrt{1-\sin^2 C}=\frac{3\sqrt{39}}{26}\).
因此
\[
\sin(B-C)
=\sin B\cos C-\cos B\sin C
=\frac{\sqrt{13}}{13}\cdot \frac{3\sqrt{39}}{26}-\frac{2\sqrt{39}}{13}\cdot \frac{5\sqrt{13}}{26}
=-\frac{7\sqrt3}{26}
\]
\end{solution}

\begin{problem}
三棱台 \(ABC-A_1B_1C_1\) 中，若 \(A_1A\perp \text{面 }ABC\)，\(AB\perp AC\)，\(AB=AC=A_1A=2\)，\(A_1C_1=1\)，\(M,N\) 分别是 \(BC,BA\) 中点.

\begin{enumerate}
\item 求证：\(A_1N\parallel \text{平面 }C_1MA\)；
\item 求平面 \(C_1MA\) 与平面 \(ACC_1A_1\) 所成夹角的余弦值；
\item 求点 \(C\) 到平面 \(C_1MA\) 的距离.
\end{enumerate}

\bitmapfigure[max width=0.52\linewidth,max totalheight=4.8cm]{img/2023/tianjin/q17_fig1.png}

\end{problem}

\begin{solution}
（1）连接 \(MN,C_1A\). 由 \(M,N\) 分别是 \(BC,BA\) 中点，得
\(MN\parallel AC\)，\(MN=\frac{AC}{2}=1\).
又由棱台性质，
\(A_1C_1\parallel AC\)，\(A_1C_1=1\).
故 \(MN\parallel A_1C_1\)，且 \(MN=A_1C_1\). 于是四边形 \(MNA_1C_1\) 为平行四边形，从而 \(A_1N\parallel MC_1\). 因为 \(MC_1\subset \text{平面 }C_1MA\)，且 \(A_1N\not\subset \text{平面 }C_1MA\)，所以 \(A_1N\parallel \text{平面 }C_1MA\).

\bitmapfigure[max width=0.52\linewidth,max totalheight=4.8cm]{img/2023/tianjin/q17_solution_fig1.png}

（2）过 \(M\) 作 \(ME\perp AC\)，垂足为 \(E\)；过 \(E\) 作 \(EF\perp AC_1\)，垂足为 \(F\)，连接 \(MF,C_1E\).

\bitmapfigure[max width=0.52\linewidth,max totalheight=4.8cm]{img/2023/tianjin/q17_solution_fig2.png}

因为 \(A_1A\perp \text{面 }ABC\)，且 \(ME\subset \text{面 }ABC\)，故 \(A_1A\perp ME\). 又 \(ME\perp AC\)，\(AC\cap A_1A=A\)，\(AC,A_1A\subset \text{平面 }ACC_1A_1\)，从而
\(ME\perp \text{平面 }ACC_1A_1\).
又 \(AC_1\subset \text{平面 }ACC_1A_1\)，故 \(ME\perp AC_1\). 又 \(EF\perp AC_1\)，\(ME\cap EF=E\)，\(ME,EF\subset \text{平面 }MEF\)，于是 \(AC_1\perp \text{平面 }MEF\).

由 \(MF\subset \text{平面 }MEF\)，故 \(AC_1\perp MF\). 于是平面 \(C_1MA\) 与平面 \(ACC_1A_1\) 所成角即 \(\angle MFE\).

又
\(ME=\frac{AB}{2}=1\)，
且 \(\cos \angle CAC_1=\frac{1}{\sqrt5}\)，\(\sin \angle CAC_1=\frac{2}{\sqrt5}\).
故
\(EF=ME\sin \angle CAC_1=\frac{2}{\sqrt5}\).
在直角三角形 \(MEF\) 中，
\(MF=\sqrt{ME^2+EF^2}=\sqrt{1+\frac{4}{5}}=\frac{3}{\sqrt5}\).
于是
\(\cos \angle MFE=\frac{EF}{MF}=\frac{2}{3}\).

（3）方法一：几何法

过 \(C_1\) 作 \(C_1P\perp AC\)，垂足为 \(P\)，作 \(C_1Q\perp AM\)，垂足为 \(Q\)，连接 \(PQ,PM\)，过 \(P\) 作 \(PR\perp C_1Q\)，垂足为 \(R\).

\bitmapfigure[max width=0.52\linewidth,max totalheight=4.8cm]{img/2023/tianjin/q17_solution_fig3.png}

由题干数据可得，\(C_1A=C_1C=\sqrt5\)，\(C_1M=\sqrt{C_1P^2+PM^2}=\sqrt5\)，根据勾股定理，\(C_1Q=\sqrt{5-\left( \frac{\sqrt2}{2} \right)^2}=\frac{3\sqrt2}{2}\).

由 \(C_1P\perp \text{平面 }AMC\)，\(AM\subset \text{平面 }AMC\)，则 \(C_1P\perp AM\)，又 \(C_1Q\perp AM\)，\(C_1Q\cap C_1P=C_1\)，\(C_1Q,C_1P\subset \text{平面 }C_1PQ\)，于是 \(AM\perp \text{平面 }C_1PQ\).

又 \(PR\subset \text{平面 }C_1PQ\)，则 \(PR\perp AM\)，又 \(PR\perp C_1Q\)，\(C_1Q\cap AM=Q\)，\(C_1Q,AM\subset \text{平面 }C_1MA\)，故 \(PR\perp \text{平面 }C_1MA\).

在 \(\mathrm{Rt}\triangle C_1PQ\) 中，\(PR=\frac{PC_1\cdot PQ}{QC_1}=\frac{2\cdot \frac{\sqrt2}{2}}{\frac{3\sqrt2}{2}}=\frac{2}{3}\).

又 \(CA=2PA\)，故点 \(C\) 到平面 \(C_1MA\) 的距离是点 \(P\) 到平面 \(C_1MA\) 的距离的两倍，\(d\bigl(C,\text{平面 }C_1MA\bigr)=2PR=\frac{4}{3}\).

方法二：等体积法

\bitmapfigure[max width=0.52\linewidth,max totalheight=4.8cm]{img/2023/tianjin/q17_solution_fig3.png}

辅助线同方法一.

设点 \(C\) 到平面 \(C_1MA\) 的距离为 \(h\).

\[
V_{C_1-AMC}
=\frac{1}{3}\times C_1P\times S_{\triangle AMC}
=\frac{1}{3}\times 2\times \frac{1}{2}\times \left( \sqrt2 \right)^2
=\frac{2}{3}
\]
\[
V_{C-C_1MA}
=\frac{1}{3}\times h\times S_{\triangle AMC_1}
=\frac{1}{3}\times h\times \frac{1}{2}\times \sqrt2\times \frac{3\sqrt2}{2}
=\frac{h}{2}
\]

由 \(V_{C_1-AMC}=V_{C-C_1MA}\)，得
\(\frac{h}{2}=\frac{2}{3}\)，
即
\(h=\frac{4}{3}\).
\end{solution}

\begin{problem}
设椭圆 \(\frac{x^2}{a^2}+\frac{y^2}{b^2}=1\)（\(a>b>0\)） 的左、右顶点分别为 \(A_1,A_2\)，右焦点为 \(F\)，已知 \(|A_1F|=3\)，\(|A_2F|=1\).

\begin{enumerate}
\item 求椭圆方程及其离心率；
\item 已知点 \(P\) 是椭圆上一动点（不与端点重合），直线 \(A_2P\) 交 \(y\) 轴于点 \(Q\)，若三角形 \(A_1PQ\) 的面积是三角形 \(A_2FP\) 面积的二倍，求直线 \(A_2P\) 的方程.
\end{enumerate}

\end{problem}

\begin{solution}
（1）由 \(|A_1F|=a+c=3\)，\(|A_2F|=a-c=1\) 解得
\(a=2\)，\(c=1\).
故 \(b=\sqrt{a^2-c^2}=\sqrt{4-1}=\sqrt3\). 所以
\(\frac{x^2}{4}+\frac{y^2}{3}=1\)，\(e=\frac{c}{a}=\frac{1}{2}\).

（2）由（1）知 \(A_2(2,0)\). 设直线 \(A_2P\) 方程为 \(y=k(x-2)\). 与椭圆联立
\[
\begin{cases}
\dfrac{x^2}{4}+\dfrac{y^2}{3}=1,\\
y=k(x-2).
\end{cases}
\]
消去 \(y\) 得
\[
(3+4k^2)x^2-16k^2x+16k^2-12=0
\]
由韦达定理并利用 \(x_{A_2}=2\) 可得
\(x_P=\frac{8k^2-6}{3+4k^2}\)，\(y_P=-\frac{12k}{3+4k^2}\)，\(Q(0,-2k)\).

\bitmapfigure[max width=0.52\linewidth,max totalheight=4.8cm]{img/2023/tianjin/q18_solution_fig1.png}

面积条件给出
\[
S_{\triangle A_2QA_1}=S_{\triangle A_1PQ}+S_{\triangle A_1A_2P}=2S_{\triangle A_2FP}+S_{\triangle A_1A_2P}
\]
化简为
\(2|y_Q|=3|y_P|\).
即
\(4|k|=3\left| \frac{12k}{3+4k^2} \right|\).
解得 \(k=\pm \frac{\sqrt6}{2}\). 故直线 \(A_2P\) 的方程为
\[
y=\pm \frac{\sqrt6}{2}(x-2)
\]
\end{solution}

\begin{problem}
已知 \(\{a_n\}\) 是等差数列，\(a_2+a_5=16\)，\(a_5-a_3=4\).

\begin{enumerate}
\item 求 \(\{a_n\}\) 的通项公式和 \(\displaystyle\sum_{i=2^{n-1}}^{2^n-1}a_i\)；
\item 已知 \(\{b_n\}\) 为等比数列，对于任意 \(k\in \N^*\)，若 \(2^{k-1}\le n\le 2^k-1\)，则 \(b_k<a_n<b_{k+1}\).

    \begin{enumerate}
    \item 当 \(k\ge 2\) 时，求证：\(2^k-1<b_k<2^k+1\)；
    \item 求 \(\{b_n\}\) 的通项公式及其前 \(n\) 项和.
    \end{enumerate}
\end{enumerate}

\end{problem}

\begin{solution}
（1）设等差数列首项为 \(a_1\)，公差为 \(d\). 由题意 \(\begin{cases} a_2+a_5=2a_1+5d=16,\\ a_5-a_3=2d=4, \end{cases}\) 解得
\(a_1=3\)，\(d=2\).
故
\(a_n=a_1+(n-1)d=2n+1\).

于是 \(\sum_{i=2^{n-1}}^{2^n-1}a_i=\sum_{i=2^{n-1}}^{2^n-1}(2i+1)=3\cdot 4^{n-1}\).

（2）（i）因为当 \(2^{k-1}\le n\le 2^k-1\) 时，
\(b_k<a_n\)，
取 \(n=2^{k-1}\)，得
\(b_k<a_{2^{k-1}}=2^k+1\).

又当 \(2^{k-2}\le n\le 2^{k-1}-1\) 时，
\(a_n<b_k\).
取 \(n=2^{k-1}-1\)，得
\(a_{2^{k-1}-1}=2^k-1<b_k\).

故
\[
2^k-1<b_k<2^k+1
\]

（2）（ii）由（1）可知
\(2^k-1<b_k<2^k+1\)，\(2^{k+1}-1<b_{k+1}<2^{k+1}+1\).
则数列 \(\{b_n\}\) 的公比 \(q\) 满足
\[
\frac{2^{k+1}-1}{2^k+1}
=2-\frac{3}{2^k+1}
<q
=\frac{b_{k+1}}{b_k}
<\frac{2^{k+1}+1}{2^k-1}
=2+\frac{3}{2^k-1}
\]
当 \(k\in \N^*\)，\(k\to +\infty\) 时，
\(2-\frac{3}{2^k+1}\to 2\)，\(2+\frac{3}{2^k-1}\to 2\)，
所以 \(q=2\).

所以
\(2^k-1<b_12^{k-1}<2^k+1\)，
即
\(2-\frac{1}{2^{k-1}}<b_1<2+\frac{1}{2^{k-1}}\).
当 \(k\in \N^*\)，\(k\to +\infty\) 时，
\(2-\frac{1}{2^{k-1}}\to 2\)，\(2+\frac{1}{2^{k-1}}\to 2\)，
所以 \(b_1=2\). 故
\(b_n=2^n\).
其前 \(n\) 项和为
\(S_n=\frac{2(1-2^n)}{1-2}=2^{n+1}-2\).
\end{solution}

\begin{problem}
已知函数 \(f(x)=\left( \frac{1}{x}+\frac{1}{2} \right)\ln(x+1)\).

\begin{enumerate}
\item 求曲线 \(y=f(x)\) 在 \(x=2\) 处切线的斜率；
\item 当 \(x>0\) 时，证明：\(f(x)>1\)；
\item 证明：
\(\frac{5}{6}<\ln(n!)-\left( n+\frac{1}{2} \right)\ln n+n\le 1\)
\end{enumerate}

\end{problem}

\begin{solution}
（1）由
\(f(x)=\frac{\ln(x+1)}{x}+\frac{1}{2}\ln(x+1)\)，
得
\(f'(x)=\frac{1}{x(x+1)}+\frac{1}{2(x+1)}-\frac{\ln(x+1)}{x^2}\).
故
\(f'(2)=\frac{1}{6}+\frac{1}{6}-\frac{\ln 3}{4}=\frac{1}{3}-\frac{\ln 3}{4}\).

（2）令
\(g(x)=\ln(x+1)-\frac{2x}{x+2}\)（\(x>0\)）.
则 \(g'(x)=\frac{1}{x+1}-\frac{4}{(x+2)^2}=\frac{x^2}{(x+1)(x+2)^2}>0\).
故 \(g(x)\) 在 \((0,+\infty)\) 上递增. 又
\(g(0)=0\)，
所以当 \(x>0\) 时，
\(g(x)>0\)，
即
\(\ln(x+1)>\frac{2x}{x+2}\).
两边乘以 \(\dfrac{1}{x}+\dfrac{1}{2}=\dfrac{x+2}{2x}>0\)，得
\(f(x)=\left( \frac{1}{x}+\frac{1}{2} \right)\ln(x+1)>1\).

（3）设
\(h(n)=\ln(n!)-\left( n+\frac{1}{2} \right)\ln n+n\).
则
\(h(n+1)-h(n)=1-\left( n+\frac{1}{2} \right)\ln\left( 1+\frac{1}{n} \right)\).
由（2）知：\(x=\dfrac{1}{n}\in (0,1]\)，则
\(f\left( \frac{1}{n} \right)=\left( n+\frac{1}{2} \right)\ln\left( 1+\frac{1}{n} \right)>1\)，
故
\(h(n+1)-h(n)<0\).
故 \(h(n)\) 在 \(n\in \N^*\) 上递减，而
\(h(1)=1\)，
故
\(h(n)\le 1\).

下证
\(\ln(n!)-\left( n+\frac{1}{2} \right)\ln n+n>\frac{5}{6}\).

令
\(\varphi(x)=\ln x-\frac{(x+5)(x-1)}{4x+2}\)（\(x>0\)）.
则
\(\varphi'(x)=\frac{(x-1)^2(1-x)}{x(2x+1)^2}\).

当 \(0<x<1\) 时 \(\varphi'(x)>0\)，\(\varphi(x)\) 递增，当 \(x>1\) 时 \(\varphi'(x)<0\)，\(\varphi(x)\) 递减，

所以 \(\varphi(x)\le \varphi(1)=0\)，故在 \(x\in (0,+\infty)\) 上
\(\ln x\le \frac{(x+5)(x-1)}{4x+2}\)
恒成立，
\[
\begin{aligned}
h(n)-h(n+1)
&=\left( n+\frac{1}{2} \right)\ln\left( 1+\frac{1}{n} \right)-1 \\
&\le \left( n+\frac{1}{2} \right)\cdot \frac{\left( 6+\frac{1}{n} \right)\frac{1}{n}}{2\left( 3+\frac{2}{n} \right)}-1 \\
&=\frac{1}{4n(3n+2)} \\
&<\frac{1}{12}\left( \frac{1}{n-1}-\frac{1}{n} \right)
\end{aligned}
\]

所以
\[
\begin{aligned}
h(2)-h(3)&<\frac{1}{12}\left( 1-\frac{1}{2} \right) \\
h(3)-h(4)&<\frac{1}{12}\left( \frac{1}{2}-\frac{1}{3} \right) \\
&\cdots \\
h(n-1)-h(n)&<\frac{1}{12}\left( \frac{1}{n-2}-\frac{1}{n-1} \right)
\end{aligned}
\]

累加得：
\(h(2)-h(n)<\frac{1}{12}\left( 1-\frac{1}{n-1} \right)\)，
而
\(h(2)=2-\frac{3}{2}\ln 2\)，\(h(1)=1\).

因为
\(\frac{7}{9}>\frac{3}{4}>\ln 2\)，
所以
\(h(2)=2-\frac{3}{2}\ln 2>\frac{5}{6}\).

则
\(-h(n)<\frac{1}{12}\left( 1-\frac{1}{n-1} \right)-2+\frac{3}{2}\ln 2\)（\(n\ge 3\)），
所以
\[
\begin{aligned}
h(1)-h(n)
&<\frac{3}{2}\ln 2-1+\frac{1}{12}\left( 1-\frac{1}{n-1} \right) \\
&<\frac{3}{2}\ln 2-1+\frac{1}{12}<\frac{1}{6}
\end{aligned}
\]
故
\(h(n)>\frac{5}{6}\)（\(n\ge 3\)）.

综上，
\(\frac{5}{6}<h(n)\le 1\).
记
\[
\begin{aligned}
u_n
&=\ln(n!)-\left( n+\frac{1}{2} \right)\ln n \\
&\quad +n
\end{aligned}
\]
则 \(u_n\in \left( \frac{5}{6},1 \right]\).
\end{solution}
